\begin{proof}
Fix \(d=[\varphi]_{\mathcal I}\), and let
\[
 \mathcal F_d:=\left\{\psi\in\PSH(X,\theta):
             \psi\leq0,\ [\psi]_{\mathcal I}=d\right\}.
\]
The family is nonempty because
\(\varphi-\sup_X\varphi\in\mathcal F_d\).  It is directed under finite
maxima: divisorial Lelong numbers satisfy
\[
 \nu(\max\{\psi_1,\psi_2\},E)
   =\min\{\nu(\psi_1,E),\nu(\psi_2,E)\},
\]
so the maximum of two members has the same divisorial profile \(d\).

Put \(U=(\sup_{\psi\in\mathcal F_d}\psi)^*\).  Choquet reduction and the
increasing-supremum formula give
\[
 \nu(U,E)=\inf_{\psi\in\mathcal F_d}\nu(\psi,E)=\nu(d,E)
\]
for every prime divisor \(E\) over \(X\).  Strong openness and the valuative
characterization of multiplier ideals therefore give
\([U]_{\mathcal I}=d\).  Since \(U\leq0\), it follows that
\(U\in\mathcal F_d\), so \(U\) is the largest member of the family.  The
family depends only on \(d\); hence \(U=\mathcal R_\theta(d)\), applying the
same construction to \(U\) returns \(U\), and the map is well defined.

The equality \(\sup_XU=0\) follows because otherwise a sufficiently small
positive constant could be added without changing the multiplier ideals,
contradicting maximality in \(\mathcal F_d\).  If \(u\) is any normalized
\(\mathcal I\)-model potential in the class \(d\), then its defining
competitor family is the same \(\mathcal F_d\), and fixedness gives
\(u=U\).  This proves the first assertion.

For the second assertion, take a compact polar generalized Cantor set
\(F\subset\mathbf P^1\) carrying a nonatomic probability measure \(\mu\).
For example, use stage lengths
\(\ell_m=\mathrm e^{-4^m}\), retain two intervals at every stage, and give
each of the \(2^m\) stage-\(m\) intervals mass \(2^{-m}\).  The masses of
nested intervals tend to zero, so \(\mu\) has no atoms; the standard
capacity criterion applies because
\[
 \sum_{m\geq0}2^{-m}\log\frac1{\ell_m}
 =\sum_{m\geq0}2^m=+\infty,
\]
and hence \(F\) is polar.

Let \(G_p\) be the normalized Green potential with
\[
 \omega+\ddc G_p=\delta_p,\qquad \sup_XG_p=0,
\]
and put
\[
 G_\mu(z):=\int_XG_p(z)\,\mathrm d\mu(p).
\]
The function is not identically \(-\infty\): away from \(F\), a positive
lower bound for the chordal distance to \(F\) gives a finite uniform lower
bound for all integrands.  Fubini's theorem yields
\[
 \omega+\ddc G_\mu=\mu.
\]
For \(0\leq s\leq1\), set \(\varphi_s=sG_\mu\).  Then
\begin{equation}
 \omega+\ddc\varphi_s=(1-s)\omega+s\mu.
 \label{eq:p148-cantor-current}
\end{equation}
On a Riemann surface, the Lelong number of a positive current at a point
is its atomic mass.  Both \(\omega\) and \(\mu\) are nonatomic, so every
Lelong number of every \(\varphi_s\) vanishes.  The valuative
characterization gives
\[
 [\varphi_s]_{\mathcal I}=[0]_{\mathcal I}.
\]

In bidegree one, the nonpluripolar product agrees with the curvature measure
on every open set where the potential is locally bounded, and it charges no
pluripolar set.  The function \(G_\mu\) is locally bounded on
\(X\setminus F\), while \(F\) is polar.  Hence
\[
 \one_{X\setminus F}
 \left\langle\omega+\ddc\varphi_s\right\rangle
 =\one_{X\setminus F}\bigl((1-s)\omega+s\mu\bigr)
 =(1-s)\one_{X\setminus F}\omega.
\]
The measure on the left gives \(F\) mass zero.  The smooth measure
\(\omega\) also gives \(F\) mass zero, since every polar subset of a Riemann
surface has zero area.  Therefore
\[
 \left\langle\omega+\ddc\varphi_s\right\rangle=(1-s)\omega,
 \qquad
 \int_X\left\langle\omega+\ddc\varphi_s\right\rangle=1-s.
\]
This is \eqref{eq:p148-same-data-all-masses}.  The first assertion applied
to their common class gives
\[
 \mathcal R_\omega([\varphi_s]_{\mathcal I})
 =\mathcal R_\omega([0]_{\mathcal I})=0.
\]
At \(s=1\), the input representative has mass zero and its canonical
normalized \(\mathcal I\)-model has mass one, proving the claimed
non-reconstruction.
\end{proof}
